% 参考文献。正文中使用 \cite{...} 引用。

\begin{thebibliography}{99}
\bibitem{tsn}
Wang L, Xiong Y, Wang Z, et al. Temporal segment networks: Towards good
practices for deep action recognition[C]//European Conference on Computer
Vision. 2016: 20--36.

\bibitem{i3d}
Carreira J, Zisserman A. Quo vadis, action recognition? A new model and the
Kinetics dataset[C]//IEEE Conference on Computer Vision and Pattern
Recognition. 2017: 6299--6308.

\bibitem{ucf101}
Soomro K, Zamir A R, Shah M. UCF101: A dataset of 101 human action classes
from videos in the wild[R]. CRCV-TR-12-01, 2012.

\bibitem{r21d}
Tran D, Wang H, Torresani L, et al. A closer look at spatiotemporal
convolutions for action recognition[C]//IEEE Conference on Computer Vision
and Pattern Recognition. 2018: 6450--6459.

\bibitem{tsm}
Lin J, Gan C, Han S. TSM: Temporal shift module for efficient video
understanding[C]//IEEE International Conference on Computer Vision. 2019:
7083--7093.

\bibitem{kinetics}
Kay W, Carreira J, Simonyan K, et al. The Kinetics human action video
dataset[J]. arXiv preprint arXiv:1705.06950, 2017.

\bibitem{imagenet}
Deng J, Dong W, Socher R, et al. ImageNet: A large-scale hierarchical image
database[C]//IEEE Conference on Computer Vision and Pattern Recognition.
2009: 248--255.

\bibitem{amp}
Micikevicius P, Narang S, Alben J, et al. Mixed precision
training[C]//International Conference on Learning Representations. 2018.

\bibitem{checkpoint}
Chen T, Xu B, Zhang C, et al. Training deep nets with sublinear memory
cost[J]. arXiv preprint arXiv:1604.06174, 2016.

\bibitem{adamw}
Loshchilov I, Hutter F. Decoupled weight decay
regularization[C]//International Conference on Learning Representations. 2019.
\end{thebibliography}
